\VerdierPhysicalPageAnchor{0042}%
\VerdierSemanticAnchor{ch1-1-3-7-cont}
exclusivement dans ce travail le langage des $G^\circ$-catégories.
Cependant, nous n'oublierons pas totalement le point de vue
\textquotedblleft catégories graduées\textquotedblright\ en utilisant
la terminologie suivante :

\VerdierSemanticAnchor{ch1-1-3-8}
\noindent\textbf{Définition 1.3.8.}\quad
Soient $G$ un groupe et $(\mathcal C,G^\circ,T,c)$ une
$G^\circ$-catégorie. Soient $g$ et $h$ deux éléments de $G$ et $X$ et
$Y$ deux objets de $\mathcal C$. Un morphisme de $\mathcal C$ :

\VerdierSemanticAnchor{ch1-1-3-8-morphism}
\[
  u\colon T(h)X\longrightarrow T(gh)Y
\]
sera appelé \emph{un morphisme de degré $g$ de $X$ dans $Y$}.

Montrons que cette terminologie n'apporte pas de confusion. En effet,
pour tout couple $h,h'$ d'éléments de $G$, le foncteur
$T(h^{-1}h')$ et les isomorphismes $c(h,h^{-1}h')$ et
$c(gh,h^{-1}h')$ définissent un isomorphisme bifonctoriel :

\VerdierSemanticAnchor{ch1-1-3-8-eq1}
\begin{equation}
  i_g(h,h')\colon
  \operatorname{Hom}_{\mathcal C}\bigl(T(h)X,T(gh)Y\bigr)
  \overset{\sim}{\longrightarrow}
  \operatorname{Hom}_{\mathcal C}\bigl(T(h')X,T(gh')Y\bigr).
  \tag{1.3.8.1}
\end{equation}

Et les identités de cocycle (1.2.1.3) montrent qu'on a la relation :

\VerdierSemanticAnchor{ch1-1-3-8-eq2}
\begin{equation}
  i_g(h,h'')=i_g(h',h'')i_g(h,h'),
  \tag{1.3.8.2}
\end{equation}
pour tout triplet $h,h',h''\in G$. Si donc, on identifie les différents
ensembles :

\VerdierSemanticAnchor{ch1-1-3-8-identification}
\[
  \operatorname{Hom}_{\mathcal C}\bigl(T(h)X,T(gh)Y\bigr)
\]
à l'aide des isomorphismes $i_g(h,h')$, la formule (1.3.8.2) montre que
ces différentes identifications sont compatibles. De plus, pour tout
quadruplet $g,g',h,h'\in G$ et tout triplet $X,Y,Z$ d'objets de
$\mathcal C$, le diagramme ci-après (1.3.8.3) est commutatif :

\VerdierSemanticAnchor{ch1-1-3-8-eq3}
\begin{center}
  \resizebox{0.92\textwidth}{!}{%
    $\begin{tikzcd}[ampersand replacement=\&,column sep=5em,row sep=4em]
      \operatorname{Hom}_{\mathcal C}\bigl(T(h)X,T(gh)Y\bigr)
      \times
      \operatorname{Hom}_{\mathcal C}\bigl(T(gh)Y,T(g'gh)Z\bigr)
      \arrow[r]
      \arrow[d,
        "{i_g(h,h')\times i_{g'}(gh,gh')}"',
        "\sim"]
      \&
      \operatorname{Hom}_{\mathcal C}\bigl(T(h)X,T(g'gh)Z\bigr)
      \arrow[d,"{i_{g'g}(h,h')}"',"\sim"]
      \\
      \operatorname{Hom}_{\mathcal C}\bigl(T(h')X,T(gh')Y\bigr)
      \times
      \operatorname{Hom}_{\mathcal C}\bigl(T(gh')Y,T(g'gh')Z\bigr)
      \arrow[r]
      \&
      \operatorname{Hom}_{\mathcal C}\bigl(T(h')X,T(g'gh')Z\bigr)
    \end{tikzcd}$%
  }
  \hfill\textup{(1.3.8.3)}
\end{center}

\VerdierSemanticAnchor{ch1-1-3-8-composition}
(les flèches horizontales sont définies par la composition dans
$\mathcal C$). Le diagramme (1.3.8.3) nous permet donc de définir sans
ambiguïté, ou plus précisément aux identifications définies par les
isomorphismes $i_g(h,h')$ près, le composé de deux morphismes de degré
respectif $g$ et $g'$ et on vérifie immédiatement que cette définition
de la composition est compatible avec la définition
